% !TEX program = pdflatex
% !BIB program = biber
% =============================================================================
%  Research article skeleton (academic-latex skill) — for preprints, course
%  papers, reports and venues without a template. When a journal or conference
%  provides a class (IEEEtran, acmart, elsarticle, llncs, revtex, ...), use that
%  instead and do not change its fonts, margins or spacing.
%  Build:  latexmk          (reads .latexmkrc; runs biber and every pass needed)
% =============================================================================
%% <<academic-latex:class>>
\documentclass[11pt, a4paper]{article}
%% <</academic-latex:class>>

%% <<academic-latex:engine>>
\usepackage[T1]{fontenc}
\usepackage{lmodern}
%% <</academic-latex:engine>>
%% <<academic-latex:language>>
\usepackage[english]{babel}
%% <</academic-latex:language>>
\usepackage{microtype}
%% <<academic-latex:geometry>>
\usepackage[margin=25mm]{geometry}
%% <</academic-latex:geometry>>

\usepackage{graphicx}
\graphicspath{{figures/}}
\usepackage{booktabs}
\usepackage{siunitx}
\usepackage{amsmath, amssymb}
\usepackage[font=small, labelfont=bf]{caption}
\usepackage{subcaption}
\usepackage{csquotes}
\usepackage{authblk}          % \author[1]{...} \affil[1]{...}

%% <<academic-latex:bibliography>>
\usepackage[backend=biber, style=numeric-comp, sorting=none]{biblatex}
%% <</academic-latex:bibliography>>
\addbibresource{references.bib}

\usepackage[hidelinks]{hyperref}
\usepackage[capitalise]{cleveref}  % always the last package

% TODO: title, authors, affiliations, corresponding-author e-mail.
\title{Informative Title That States the Main Finding}
\author[1]{First Author\thanks{Corresponding author: \href{mailto:first.author@example.org}{first.author@example.org}}}
\author[2]{Second Author}
\affil[1]{Department, University, City, Country}
\affil[2]{Institute, Organisation, City, Country}
\date{}                       % or \today for a dated preprint
\hypersetup{pdftitle={Informative Title That States the Main Finding}, pdfauthor={First Author, Second Author}}

\begin{document}
\maketitle

\begin{abstract}
% TODO: 150-250 words unless the venue says otherwise. Context, gap, aim,
% method, main results with numbers, implication. No citations.
Context \ldots{} However, \ldots{} Here we \ldots{} We find that \ldots{} These
results suggest \ldots
\end{abstract}

\noindent\textbf{Keywords:} first keyword; second keyword; third keyword

\section{Introduction}\label{sec:introduction}
% Context -> gap -> what this paper does -> contributions -> outline.
% TODO: write the introduction.
Typesetting with \LaTeX{} \autocite{lamport1994latex} separates content from
presentation; literate programming \autocite{knuth1984literate} applies a
related idea to code.

\section{Related work}\label{sec:related-work}
% TODO: synthesise prior work by theme and position this paper.

\section{Methods}\label{sec:methods}
% TODO: enough detail to reproduce the work (data, tools with versions, parameters).
\Cref{fig:overview} gives an overview.

\begin{figure}[htbp]
  \centering
  \includegraphics[width=0.6\linewidth]{example-image}  % placeholder (mwe package)
  \caption{Overview of the approach.}
  \label{fig:overview}
\end{figure}

\section{Results}\label{sec:results}
% TODO: report results; call out every figure and table before it appears.
\Cref{tab:results} summarises the results.

\begin{table}[htbp]
  \centering
  \caption{Main results (mean over five runs).}
  \label{tab:results}
  \begin{tabular}{l S[table-format=2.1] S[table-format=3.0]}
    \toprule
    {Method} & {Accuracy (\%)} & {Time (\unit{\milli\second})} \\
    \midrule
    Baseline & 71.2 & 120 \\
    Proposed & 78.9 & 135 \\
    \bottomrule
  \end{tabular}
\end{table}

\section{Discussion}\label{sec:discussion}
% TODO: interpretation, comparison with prior work, limitations.

\section{Conclusion}\label{sec:conclusion}
% TODO: the answer to the research question and its implications; no new results.

% Statements many venues require; delete the ones yours does not ask for.
\section*{Acknowledgements}
% TODO: people and funding (agency, grant number).

\section*{Data availability}
% TODO: where the data and code can be obtained (repository + DOI), or why not.

\section*{Author contributions}
% TODO: e.g. CRediT roles: Conceptualization, Methodology, Software, Writing -- original draft.

\printbibliography

\appendix
\section{Supplementary details}\label{app:details}
% TODO: add material or delete the appendix.

\end{document}
